\section{Formulation and research programme}\label{sec:program}

Throughout this appendix, a \emph{rectification} is an automorphism
of the entire affine space that sends the given parametrized curve
to the standard line. The choice of field and the requirement of a
polynomial inverse are part of the problem. Birational rectification
is a weaker condition.

\subsection{Lifting the parameter to a coordinate}

Let $k$ be a field, let $R=k[x,y,z]$, and write
\[
 \gamma(t)=(p(t),q(t),r(t)),\qquad p,q,r\in k[t].
\]
The parametrization defines a closed embedding
$\A^1_k\hookrightarrow\A^3_k$ if and only if
\begin{equation}\label{eq:closed-embedding}
 k[p,q,r]=k[t].
\end{equation}
This is equivalent to surjectivity of the restriction homomorphism
$\gamma^*:R\to k[t]$. In particular, there is a polynomial $H_0\in R$
with $H_0\gamma=t$. Such a polynomial need not, however, be a
\emph{coordinate}, that is, a component of a polynomial automorphism.

\begin{lemma}[Coordinate lift]\label{lem:coordinate-lift}
The embedding $\gamma$ admits a rectifier $\Phi\in\Aut_k(\A^3)$
with $\Phi\gamma(t)=(t,0,0)$ if and only if there is a coordinate
$H\in R$ satisfying $H\gamma=t$.
\end{lemma}
\begin{proof}
Necessity follows from the first component of $\Phi$. Conversely,
complete $H$ to an automorphism $(H,U,V)$. Its value on the curve
is $(t,u(t),v(t))$. Follow it by the automorphism
\[
 (T,Y,Z)\longmapsto(T,Y-u(T),Z-v(T)).
\]
The resulting composition rectifies $\gamma$.
\end{proof}
This criterion is known and is used explicitly in
\cite[Lemma~1]{Lewis2024}. Its constructive content is the choice
of a coordinate among all lifts $H_0+\ker\gamma^*$.

\subsection{What the two-dimensional case provides}

In characteristic zero, the Abhyankar--Moh--Suzuki theorem states
that every closed embedding $\A^1_k\hookrightarrow\A^2_k$ is
rectifiable \cite{AM1975,Suzuki1974}. For unparametrized images,
this means that the action of $G_2(k)=\Aut_k(\A^2)$ is transitive.
Setting $L_2=\{y=0\}$ gives the following description of the set
of such images:
\[
 \{\text{images of closed embeddings }\A^1_k\hookrightarrow\A^2_k\}
 \simeq G_2(k)/\operatorname{Stab}(L_2).
\]
Here the quotient consists of left cosets
$g\operatorname{Stab}(L_2)$, with
\begin{equation}\label{eq:plane-stabilizer}
 \operatorname{Stab}(L_2)=
 \{(x,y)\mapsto(ax+P(y),by):a,b\in k^*,\ P\in k[y]\}.
\end{equation}
Indeed, preservation of the ideal $(y)$ forces the image of $y$
to be $by$. After scaling, an automorphism of $k[y][x]$ remains,
so the image of $x$ is $ax+P(y)$ with $a\in k^*$. The reverse
inclusion is immediate.

Three classical statements should be distinguished.
The Abhyankar--Moh--Suzuki theorem describes the orbit of embedded
lines. The Jung--van der Kulk theorem describes $G_2(k)$ through
affine and triangular automorphisms and holds over every field
\cite[Theorem~1]{EV2001}. The Noether--Castelnuovo theorem states
that $\Cr_2(\C)$ is generated by projective automorphisms and the
standard quadratic involution; see the discussion in
\cite[\S1]{BH2015}. Describing the orbit of a line does not by
itself describe the whole Cremona group. In positive
characteristic, the general plane rectification theorem is false,
so its conclusion cannot be transferred to an arbitrary field
without qualification.

\subsection{Birational rectification and the precise group formulation}

Choosing the standard affine chart gives an inclusion
\[
 G=\Aut_k(\A^3)\subset B=\Cr_3(k)=\operatorname{Bir}_k(\mathbb P^3).
\]
For $b\in B$, there is an exact criterion:
\begin{equation}\label{eq:ring-preservation}
 b\in G\quad\Longleftrightarrow\quad b^*R=R.
\end{equation}
Indeed, this equality means that all components of both $b$ and
$b^{-1}$ are polynomial. In terms of orders along irreducible
polynomials $P\in R$, the condition is
\[
 v_P(b_i)\geq0,\qquad v_P((b^{-1})_i)\geq0
 \quad\text{for every }P\text{ and }i.
\]
This uses unique factorization in $R$. Poles along the hyperplane
at infinity are allowed. Checking the Jacobian alone does not
replace condition~\eqref{eq:ring-preservation}.

\begin{proposition}\label{prop:birational-any-field}
Over any field $k$, every embedding~\eqref{eq:closed-embedding}
has a birational rectifier over $k$. The rectifier and its inverse
can be chosen regular near the generic points of the respective
curves. The rational Jacobian determinant of the rectifier can be
made equal to $1$.
\end{proposition}
\begin{proof}
Differentiating $H_0(p,q,r)=t$ gives
\[
 (H_0)_x(p,q,r)p'+(H_0)_y(p,q,r)q'+(H_0)_z(p,q,r)r'=1.
\]
Thus, after cyclically permuting the coordinates, we may assume
$p'\ne0$. The extension $k(t)/k(p)$ is finite and separable, and
$k(t)=k(p)(q,r)$. By the primitive element theorem, we can choose
$\lambda(p)\in k[p]$ such that
$k(p,q+\lambda(p)r)=k(t)$. Indeed, only finitely many coefficients
$\lambda\in k(p)$ are bad, whereas $k[p]$ is infinite even when
$k$ is finite. This also follows directly by comparing the
distinct embeddings of the separable extension: each pair of
distinct embeddings excludes at most one value of $\lambda$.

Apply the shear $(x,y,z)\mapsto(x,y+\lambda(x)z,z)$ and again
denote the second component of the curve by $q$. There is a
$\tau\in k(x,y)$ satisfying $\tau(p,q)=t$; the denominator of the
chosen expression does not vanish identically on the curve. Set
\begin{equation}\label{eq:birational-construction}
 s=z-r(\tau(x,y))+\tau(x,y),\qquad
 b=(s,x-p(s),y-q(s)).
\end{equation}
Then $b\gamma=(t,0,0)$. For input $(T,U,V)$, the inverse is given by
\begin{equation}\label{eq:birational-inverse}
 \begin{split}
 X&=U+p(T),\qquad Y=V+q(T),\\
 b^{-1}(T,U,V)&=(X,Y,T+r(\tau(X,Y))-\tau(X,Y)).
 \end{split}
\end{equation}
Both compositions are the identity by direct substitution.
On the line $(T,0,0)$, the denominator becomes the same nonzero
element of $k(T)$, so the inverse is also defined at the generic
point. All shears used have Jacobian determinant $1$, as do the
cyclic permutations.
\end{proof}

Over an algebraically closed field, birational equivalence of the
projective closure of our curve to a line is a special case of the
Mella--Polastri theorem: two birational embeddings of a
$d$-dimensional variety into $\mathbb P^n$ are Cremona equivalent
when $n\geq d+2$ \cite[Theorem~1]{MP2009}. The elementary proof
above specifies a map over the original field. Neither statement
guarantees regularity at every point of affine space.

For Shastri's curve, one can explicitly take
\[
 \tau(x,y)=\frac{x(y+1)}{x^2-y-3},\qquad
 s=z-\tau^5+11\tau,\qquad
 b=(s,x-s^3+3s,y-s^4+4s^2).
\]
The identity $p(q+1)=t(p^2-q-3)$, for $p=t^3-3t$ and
$q=t^4-4t^2$, verifies recovery of the parameter. This is a
birational rectifier over $\Q$ with poles allowed; both
compositions with its explicit inverse are checked in the new
computational package. Its existence is compatible with the
impossibility of polynomial rectification over $\Q$ discussed below.

Let $L=\{y=z=0\}$. Denote by $S_L$ the subgroup of $B$ consisting
of birational maps that induce isomorphisms between some
neighbourhoods of the generic point of $L$ and act identically
on $L$. Thus the regularity requirement applies to the inverse
map as well.

\begin{proposition}[The coset-intersection problem]\label{prop:coset}
If $b$ is a birational rectifier from
Proposition~\ref{prop:birational-any-field}, then
\begin{equation}\label{eq:coset}
 \gamma\text{ is polynomially rectifiable}
 \quad\Longleftrightarrow\quad S_Lb\cap G\ne\varnothing.
\end{equation}
\end{proposition}
\begin{proof}
For a polynomial rectifier $\Phi$, the map $h=\Phi b^{-1}$
belongs to $S_L$, and $\Phi=hb$. Conversely, if $hb\in G$ with
$h\in S_L$, then $(hb)\gamma(t)=(t,0,0)$ as a rational identity.
Since both sides are polynomial, the identity holds everywhere.
\end{proof}
The original problem is therefore reformulated as the search for
a correction $h$ that removes the affine poles of $b$ and its
inverse while preserving the parametrized line.
Formula~\eqref{eq:coset} is not by itself an algorithm and does
not prove that the required intersection is nonempty.

\subsection{The relative subgroup as a working mechanism}

The stabilizer of the rational function $x$ and its polynomial
part are
\[
 J_x=\{b\in B:b^*x=x\}\simeq\Cr_2(k(x)),\qquad
 G_x=G\cap J_x=\Aut_{k[x]}k[x,y,z].
\]
An element of $G_x$ becomes a tame automorphism in two variables
over $k(x)$, but its factorization over that field may have
denominators. Eliminating them over $k[x]$ is the substantive
part of the construction. This approach has direct predecessors
in the theory of coordinate transfer and residual coordinates
\cite{EV2001,Lewis2024}.

If $A=(x,H,W)\in G_x$, $H\gamma=t$, and $W\gamma=w(t)$, then
\begin{equation}\label{eq:relative-rectifier}
 \Phi=(H,x-p(H),W-w(H))
\end{equation}
rectifies $\gamma$. Within $G_x$ itself, the natural target
parametrization is $(p(t),t,0)$: maps fixing $x$ do not change
the first component $p(t)$. A subsequent external tame
automorphism sends this parametrization to $(t,0,0)$.

\begin{remark}\label{rem:relative-not-necessary}
For a fixed choice of $x$, the existence of a relative coordinate
$H$ with $H\gamma=t$ is sufficient for rectifiability.
Its necessity for an arbitrary rectifiable embedding is not
proved here. Allowing an arbitrary preliminary polynomial change
of variables without further restrictions would make a
completeness claim circular: that change could itself be the
rectifier being sought.
\end{remark}

\subsection{The field of definition and the descent obstruction}

For Shastri's curve
\[
 \gamma_T=(t^3-3t,t^4-4t^2,t^5-10t),
\]
the real image represents a trefoil
\cite[\S3, Example~1]{Shastri1992}.
A polynomial automorphism of $\mathbb R^3$ is a proper
diffeomorphism and extends to a homeomorphism of the one-point
compactification $S^3$. It cannot send a nontrivial knot to an
unknot. Therefore rectification over $\mathbb R$, and hence over
$\Q$, is impossible. The same argument applies to any number
field with a real embedding: apply that embedding to the
coefficients of a hypothetical rectifier.

Together with the explicit rectifier over $K=\Q(\sqrt{-3})$,
this yields
\[
 \min\{[F:\Q]: F\text{ a number field},\
 \gamma_T\text{ rectifiable over }F\}=2.
\]
A suitable number field must be totally imaginary.
Uniqueness of $K$ among quadratic fields is not established;
the absence of a root of $a^2-a+1$ in a field excludes direct
specialization of our formula, but not a different construction.

Polynomial descent has a precise group formulation. Let
$\alpha\in\Aut_K(\A^3)$ rectify $\gamma_T$, let
$\Gamma=\operatorname{Gal}(K/\Q)$, and let $H_L(K)$ be the
pointwise \emph{polynomial} stabilizer of $L$. Then
\[
 c_\sigma=\alpha\,\sigma(\alpha)^{-1}\in H_L(K),\qquad
 c_{\sigma\tau}=c_\sigma\,\sigma(c_\tau).
\]
A rectifier over $\Q$ exists if and only if there is an
$h\in H_L(K)$ with $c_\sigma=h^{-1}\sigma(h)$ for every $\sigma$.
Indeed, this is equivalent to $h\alpha$ and its polynomial
inverse being defined over $\Q$. For Shastri's curve, no such
$h$ exists. If the polynomial stabilizer is replaced by $S_L$,
the obstruction disappears: for a birational rectifier $b$
over $\Q$, one can take $h=b\alpha^{-1}$.
Passing to the full Cremona group thus loses, among other
things, information about polynomial descent.
Only an elementary cocycle identity is used here; no
classification of the nonabelian set $H^1(\Gamma,H_L(K))$
is asserted.

\subsection{What remains a programme rather than a result}

\begin{enumerate}
\item \textbf{Completeness of transfers.} Describe a class of
parametrizations for which the choice of modification centres
and exceptional divisors always leads to a relative coordinate.
Finite computations for individual indices do not establish
such completeness.
\item \textbf{Controlled changes of projection.} Determine which
prescribed affine or bounded polynomial changes extend the
applicability of $G_x$ without a circular formulation.
\item \textbf{The tame orbit of the curve itself.} Wildness of
the constructed rectifier does not exclude another tame
rectifier. All coordinates in $H+\ker\gamma^*$ must be studied,
rather than only the chosen $H$.
\item \textbf{Arithmetic.} Determine other fields of rectification;
for example, test whether Shastri's curve has a rectifier over
$\Q(i)$. The absence of a root needed by a particular formula
and an obstruction to the existence of every automorphism are
different statements.
\end{enumerate}
These aims do not require a complete description of $\Cr_3(\C)$
in advance. They concern specific subgroups, cosets, and regularity
conditions on a fixed affine model.
